% =========================================================
% CHAPTER 1: INTRODUCTION
% =========================================================

\chapter{INTRODUCTION}
\label{chap:introduction}

Modern database systems are central to nearly every software application, from
e-commerce platforms to enterprise analytics pipelines. The efficiency with
which a database executes SQL queries directly affects application
responsiveness, server resource consumption, and the overall user experience.
Among the many factors that determine query efficiency, the choice of execution
plan made by the query optimizer is one of the most consequential. A
well-chosen plan can execute a query in milliseconds; a poor one can take
seconds or minutes on the same data.

Traditional database optimizers, including the optimizer in MySQL 8.0, make
these plan selection decisions using internally maintained cost models and
statistical summaries of the stored data. While this approach performs well
under stable conditions, it faces a fundamental challenge in real-world
deployments: workloads are not static. The mix of queries arriving at a
database changes over time as application usage patterns evolve, new features
are deployed, and seasonal or event-driven fluctuations occur. A plan that was
optimal for the previous workload distribution may become inefficient as the
distribution shifts.

QueryPilot addresses this challenge as an external Python middleware that wraps
MySQL query execution. For every incoming SQL query, QueryPilot extracts
structural features, generates a small set of candidate execution variants
using MySQL optimizer hints, predicts their relative performance using an
online-trained machine learning ensemble, estimates prediction uncertainty,
monitors the query stream for workload drift, and uses a guarded decision
controller called SAGE to make a final bounded-risk execution decision.

\section{OVERVIEW}
\label{sec:overview}

QueryPilot is an uncertainty-aware, self-adaptive query optimization framework
designed to operate as an external layer over MySQL 8.0 without requiring any
modifications to the database internals. Applications interact with QueryPilot
through a Python SDK interface rather than connecting directly to MySQL.
Internally, every query passes through a ten-stage pipeline that includes
workload collection, feature extraction, candidate generation, performance
prediction, uncertainty estimation, drift detection, SAGE decision-making,
query execution, feedback collection, and model update. The results of every
experiment are visualized through a real-time Streamlit dashboard.

The central research question that QueryPilot investigates is: \textit{Can an
uncertainty-aware, stability-constrained query optimization framework adapt to
dynamic database workloads while reducing query execution time and controlling
performance regressions?} The project evaluates this question experimentally
using the TPC-H benchmark at scale factor 1 under four controlled workload
scenarios: stable, sudden shift, gradual shift, and data skew.

\section{PROBLEM STATEMENT}
\label{sec:problem_statement}

The problem addressed in this project has two distinct dimensions. First,
database workloads are non-stationary: the query patterns arriving at a
production database change over time, and an execution plan that was optimal
under one distribution may become inefficient under another. MySQL's native
optimizer addresses this partially through periodic statistics collection, but
it does not explicitly monitor for workload shifts or adapt its strategy in
response to them.

Second, learned query optimization techniques --- which use machine learning to
improve execution plan selection --- suffer from a fragility problem. A model
trained under one workload distribution can make confidently wrong
recommendations when that distribution drifts. An unconstrained ML-based
optimizer that blindly follows its model's recommendations under these
conditions can introduce performance regressions, making queries slower than
MySQL's native optimizer would have. Neither the problem of workload drift
detection nor the problem of uncertainty-aware candidate filtering has been
adequately addressed in the context of external hint-based query optimization
for MySQL.

\subsection{Motivation}
\label{subsec:motivation}

The motivation for QueryPilot stems from two observed gaps in the existing
landscape of query optimization. Traditional rule-based optimizers make static
decisions based on collected statistics and do not adapt when those statistics
become unrepresentative due to workload shifts. Learned optimizers such as Bao
\cite{marcus2021bao} and Neo \cite{marcus2019neo} demonstrate that machine
learning can improve over native optimizers on training workloads, but neither
system explicitly addresses what should happen when the model's predictions are
unreliable or when the workload has drifted away from the training distribution.
QueryPilot is motivated by the hypothesis that a conservative,
uncertainty-aware decision framework --- one that defers to the native optimizer
when conditions are unfavourable --- is more robust in practice than a system
that always applies its best ML recommendation.

\section{OBJECTIVES}
\label{sec:objectives}

The objectives of this project are:

\begin{itemize}[leftmargin=2em]
  \item To design and implement QueryPilot, a Python middleware for adaptive
        query optimization over MySQL 8.0 that operates without modifying
        database internals.
  \item To build a feature extraction pipeline that derives eight meaningful
        query characteristics from MySQL's EXPLAIN output in real time.
  \item To implement a candidate optimization generator that produces
        controlled query variants using MySQL's supported optimizer hint
        system.
  \item To train a three-model ensemble machine learning predictor on observed
        execution data and use ensemble disagreement as a practical uncertainty
        estimator.
  \item To implement the Page-Hinkley sequential change-point detection
        algorithm for real-time workload drift detection.
  \item To design and implement SAGE, a guarded decision controller that
        combines uncertainty estimation, drift detection, and risk-bounded
        adaptation.
  \item To evaluate QueryPilot against a native MySQL baseline and an
        unconstrained ML-Only baseline across four controlled workload
        scenarios using the TPC-H benchmark with an interleaved execution
        methodology.
  \item To build a Streamlit monitoring dashboard that visualizes query
        performance, SAGE decisions, drift detection events, and prediction
        uncertainty in real time.
\end{itemize}

\subsection{Scope of the Project}
\label{subsec:scope}

The scope of this project is bounded to MySQL 8.0.43 as the target database
system. Optimization decisions are made using MySQL's supported optimizer hint
system, specifically the NO\_INDEX, HASH\_JOIN, and NO\_BKA hints. The
framework does not claim to enumerate or evaluate every possible execution plan.
Evaluation uses the TPC-H benchmark at scale factor 1, generating approximately
8.6 million rows across eight tables. The framework is implemented as a Python
SDK and does not implement a transparent wire-protocol database proxy. Extensions
such as a proxy layer, support for other database systems, and advanced hint
types such as FORCE INDEX or JOIN ORDER are identified as future work. All
experiments are conducted on a single machine running Windows, with MySQL and
the Python middleware co-located on the same host.
